\documentclass[11pt]{article}
\usepackage[margin=1in]{geometry}
\usepackage{amsmath,amssymb,amsthm,hyperref}
\newtheorem{lemma}{Lemma}
\newtheorem{theorem}[lemma]{Theorem}
\newcommand{\R}{\mathbb R}
\newcommand{\norm}[1]{\left\lVert #1\right\rVert}
\newcommand{\tr}{\operatorname{tr}}
\title{Charged construction of relative fractional resolvent factors\\from feasible electrical flows}
\author{Independent dynamic session: constructive subtask}
\date{November 2026}
\begin{document}
\maketitle

Let $L=B^TWB$ be a connected weighted graph Laplacian on $n\ge2$
vertices and $m$ edges.  The diagonal conductance matrix $W$ is positive.
All negative powers of $L$ are zero on the constants. Fix an edited edge
vector $u=\sqrt d\,(e_a-e_b)$, $d>0$, and write
\[
 \tau=u^TL^\dagger u>0,\qquad v=L^\dagger u,
 \qquad R_t=(L+tI)^{-1},\quad t>0.
\]
The endpoints need not be adjacent in the original graph. All statements
below concern exact-real arithmetic and explicit vector output.

\begin{lemma}[Clipping the initial potential]\label{lem:clip}
Suppose $0<\eta<1$ and an approximate potential $q$ satisfies
\[
 \norm{q-v}_L\le\eta\sqrt\tau.
\]
Clip every coordinate of $q$ to $[q_b,q_a]$ and subtract the mean of
the clipped vector, obtaining $\widehat v$. This operation is well
defined, requires $O(n)$ arithmetic operations, and satisfies
\[
 \norm{\widehat v-v}_L\le\eta\sqrt\tau,
 \qquad
 \norm{\widehat v}_2\le
       \sqrt n(1+\eta)\tau/\sqrt d,
 \qquad \widehat v\perp\mathbf1.
\]
\end{lemma}
\begin{proof}
Cauchy--Schwarz in the $L$ energy gives
\[
 |u^T(q-v)|=|v^TL(q-v)|\le\eta\tau.
\]
Since $u^Tv=\tau$, the endpoint difference $q_a-q_b$ is positive
and at most $(1+\eta)\tau/\sqrt d$.
Let $\widetilde q$ be the clipped vector. Scalar clipping is
nonexpansive, so every graph edge contribution to the Dirichlet energy
decreases: $\widetilde q^TL\widetilde q\le q^TLq$.
Both edited endpoints remain fixed, so $u^T\widetilde q=u^Tq$.
The identity
\[
 \norm{z-v}_L^2=z^TLz-2u^Tz+\tau
\]
therefore proves that clipping does not increase the energy error.
Subtracting a constant does not affect this error. After recentering,
each coordinate has absolute value at most the original clipped range,
proving the asserted Euclidean bound.
\end{proof}

The fact that the initial right-hand side is an edge vector is used
precisely in this clipping step. Clipping to its two endpoint values
requires no knowledge of the exact value of $\tau$.

\begin{lemma}[Feasible grounded flows control two norms simultaneously]
\label{lem:ground}
Fix a mean-zero vector $w\in\R^n$. Add a ground vertex and an edge of
conductance $t$ from each original vertex to ground. A feasible augmented
flow consists of an original-edge flow $f\in\R^m$ and a vector of
ground-edge flows $g\in\R^n$ such that
\[
 B^Tf+g=w.
\]
The ground demand is zero. Define its energy by
\[
 \mathcal E(f,g)=f^TW^{-1}f+t^{-1}\norm g_2^2.
\]
The minimizing flow is
\[
 y=R_tw,\qquad f_*=WBy,\qquad g_*=ty,
 \qquad \mathcal E_*=w^TR_tw.
\]
Suppose the feasible pair obeys
\[
 \mathcal E(f,g)-\mathcal E_*\le\zeta^2\mathcal E_*.
\]
Then the explicitly constructed vector $\widehat x=B^Tf=w-g$ obeys,
with $x=LR_tw$,
\begin{align*}
 \norm{\widehat x-x}_{L^\dagger}^2
       &\le\zeta^2\mathcal E_*,\\
 \norm{\widehat x-x}_2^2
       &\le\zeta^2t\mathcal E_*,\\
 \norm{L^{-\alpha/2}(\widehat x-x)}_2^2
       &\le\zeta^2\mathcal E_*t^{1-\alpha}
        \le\zeta^2\norm w_2^2t^{-\alpha},\qquad 0<\alpha<1.
\end{align*}
All these vectors are mean zero. If $0<\zeta\le1$, the same output
also has explicit witnesses
\[
 f^TW^{-1}f\le 2\norm w_2^2/t,
 \qquad \norm{\widehat x}_2\le3\norm w_2.
\]
\end{lemma}
\begin{proof}
The equations for the minimizer follow from the electrical flow
optimality conditions and $(L+tI)y=w$. Feasibility implies
\[
 e:=\widehat x-x=B^T(f-f_*)=-(g-g_*).
\]
The difference between the feasible flow and the minimizer is a
circulation in the augmented graph. Its energy inner product with the
electrical flow is zero. Consequently
\[
 \mathcal E(f,g)-\mathcal E_*
  =\norm{f-f_*}_{W^{-1}}^2+t^{-1}\norm{g-g_*}_2^2.
\]
The minimum original-graph flow energy with divergence $e$ is
$e^TL^\dagger e$. Thus each summand bounds the corresponding norm
claimed above. Spectral H\"older interpolation on $\mathbf1^\perp$
gives
\[
 e^TL^{-\alpha}e
 \le (e^TL^\dagger e)^\alpha(e^Te)^{1-\alpha}
 \le\zeta^2\mathcal E_*t^{1-\alpha}.
\]
Also $\mathcal E_*\le\norm w_2^2/t$.
Finally, $\mathcal E(f,g)\le2\mathcal E_*$ bounds the original flow
energy and implies $\norm g_2\le\sqrt2\norm w_2$. Use
$\widehat x=w-g$ for the last assertion. The ground demand and
$w\perp\mathbf1$ force $\sum g_i=0$; divergence of an original
flow also has coordinate sum zero.
\end{proof}

\begin{theorem}[Relative resolvent factor from a clipped potential]
\label{thm:factor}
Use the clipped $\widehat v$ of Lemma~\ref{lem:clip} as the demand in
Lemma~\ref{lem:ground}, and output $\widehat x_t=B^Tf_t$. Let
\[
 K_\alpha=\sup_{z>0}\frac{z^{(1-\alpha)/2}}{1+z}\le1.
\]
Then
\[
 t^{(1+\alpha)/2}
 \norm{L^{-\alpha/2}(\widehat x_t-R_tu)}_2
 \le K_\alpha\eta\sqrt\tau
             +\zeta\sqrt t\,\norm{\widehat v}_2.
\]
In particular, a common normalized column-error target $\xi>0$
over $0<t\le t_{\max}$ is achieved by
\[
 \eta\le\frac{\xi}{2\sqrt\tau},\qquad
 \zeta\le\frac{\xi\sqrt d}
 {2\sqrt{nt_{\max}}(1+\eta)\tau},
\]
subject to $\eta<1$ and $\zeta\le1$. Neither accuracy prescription
involves $\lambda_2(L)$ or a condition number.
\end{theorem}
\begin{proof}
The identity $LR_tv=R_tu$ holds on $\mathbf1^\perp$.
For $h=\widehat v-v$, spectral calculus gives
\[
 t^{(1+\alpha)/2}
 \norm{L^{-\alpha/2}LR_th}_2
 \le K_\alpha\norm h_L.
\]
Indeed the multiplier from $L^{1/2}h$ to the vector on the left is
$t^{(1+\alpha)/2}\lambda^{(1-\alpha)/2}/(\lambda+t)$.
Apply the triangle inequality and Lemma~\ref{lem:ground} with
$w=\widehat v$.
\end{proof}

\paragraph{Charged primary primitive and probability amplification.}
Kelner, Orecchia, Sidford and Allen-Zhu~\cite{KOSZ}, Theorem 7.4 and
Algorithm 6, return a feasible flow whose expected energy is at most
$(1+\rho)$ times optimum in
$O(M\log^2N\log\log N\log(1/\rho))$ arithmetic operations on an
$N$-vertex, $M$-edge weighted graph. The simpler Theorem 3.2 adds
only a $\log N$ inside the final logarithm and is sufficient here.
These bounds depend on graph size and requested accuracy, without a
weight-ratio or spectral-gap logarithm.

For Lemma~\ref{lem:ground}, run this primitive with
$\rho=\zeta^2/4$. Markov's inequality bounds the probability that
the energy excess exceeds $\zeta^2\mathcal E_*$ by $1/4$.
Run $s=O(\log(1/\delta))$ independent copies and select the feasible
output of smallest energy. Failure probability is at most $4^{-s}$;
selection requires no knowledge of $\mathcal E_*$. The augmented graph
has $n+1$ vertices and $m+n$ edges. Its complete cost is therefore
\[
 \widetilde O\bigl((m+n)\log(1/\zeta)\log(1/\delta)\bigr),
\]
where suppressed logarithms concern $n$ alone. To materialize all flow
coordinates, store the off-tree flow values updated during cycle
corrections and recover the remaining tree flows from the residual
demands by a subtree-sum pass. This costs $O(m+n)$; forming $B^Tf$
costs the same. Thus the claim includes construction of the vector
factor and its flow witness.

For the initial potential, any usual relative-energy Laplacian solver
provides $q$ at the same type of cost with logarithmic dependence on
$1/\eta$. Its high-confidence guarantee can be amplified by independent
repetition, selecting the minimum of
$q^TLq-2u^Tq$, since this objective plus $\tau$ equals the squared
energy error. The squared-error expectation needed for this selection
is also proved in the KOSZ analysis.

\paragraph{Consequence for a small compression Gram matrix.}
The original flow witnesses make a later use of $L^{-\alpha}$ charged
as well. Suppose an $n\times q$ factor $Z$ has columns
$z_j=B^Tf_j$ and all $f_j$ are explicit. Let
\[
 H=\sum_j f_j^TW^{-1}f_j,\qquad N=\norm Z_F^2.
\]
For $G=Z^TL^{-\alpha}Z$ the negative-power integral is
\[
 G=\frac{\sin(\pi\alpha)}\pi\int_0^\infty
       t^{-\alpha}Z^T(L+tI)^{-1}Z\,dt.
\]
The low tail $0<t<T$ has operator norm at most
$c_\alpha H T^{1-\alpha}/(1-\alpha)$, and the high tail $t>T$
has operator norm at most $c_\alpha N T^{-\alpha}/\alpha$.
This uses $\tr(Z^TL^\dagger Z)\le H$ and
$\tr(Z^TZ)=N$. On the retained interval ordinary shifted
energy solves suffice: columnwise relative energy accuracy $\theta$
gives operator error at most $\theta N/t$ in the symmetrized matrix
$Z^T(L+tI)^{-1}Z$. No relative inverse power is applied to these
solve errors. A logarithmic quadrature can therefore approximate the
$q\times q$ Gram matrix from $q$ shifted solves per node, with the
cutoff ratio and accuracy dependence explicitly exposed through $H,N$.
The scale dependence can be made explicit. Put $a=N/H$ and
$S=H^\alpha N^{1-\alpha}$, assuming $Z\ne0$. After $t=ae^y$, the
analytic integrand on any fixed strip of width less than $\pi$ is bounded
by
\[
 C_\alpha S\min(e^{(1-\alpha)y},e^{-\alpha y}).
\]
Consequently the usual logarithmic trapezoid construction gives absolute
Gram error $\gamma$ with
$O_\alpha(\log^2(2+S/\gamma))$ nodes and ordinary solve accuracy
of order $(\gamma/(C_\alpha S))^{1/(1-\alpha)}$.
To verify the last assertion, put $\rho=\gamma/(C_\alpha S)$.
The low cutoff is $a\rho^{1/(1-\alpha)}$; integrating the solve-error
bound $\theta N/t$ produces at most
$C_\alpha\theta S\rho^{-\alpha/(1-\alpha)}$.

For raw positive trapezoid factors of the form
\[
 z_j=\sqrt{\omega_j}\,\widehat x_{t_j},\qquad
 \omega_j\le c_\alpha h t_j^{\alpha+1},\qquad
 t_j=ae^{jh},\quad0<h\le1,
\]
Lemma~\ref{lem:ground} and geometric sums give
\[
 H\le C_\alpha\norm{\widehat v}_2^2t_{\max}^{\alpha},\qquad
 N\le C_\alpha\norm{\widehat v}_2^2t_{\max}^{\alpha+1}.
\]
It follows that
\[
 S\le C_\alpha\norm{\widehat v}_2^2t_{\max}
 \le C_\alpha n(1+\eta)^2\tau^2(t_{\max}/d).
\]
Thus a polynomial bound on $t_{\max}/d$ in the accuracy and leverage
parameters suffices for a Gram computation with no spectral-gap or
conductance-ratio logarithm. Its complete cost consists of the number
of nodes times $q$ standard shifted solves, $O(nq^2)$ arithmetic per
node for the small Gram products, and scalar quadrature operations.

\paragraph{Why ordinary shifted energy error is insufficient for the
raw factor.}
Here is an actual three-vertex graph example. For $0<\kappa<1$, use
conductances
\[
 w_{13}=w_{23}=\kappa/3,\qquad
 w_{12}=(1-\kappa/3)/2.
\]
The normalized mean-zero vectors
$p=(1,-1,0)^T/\sqrt2$ and $r=(1,1,-2)^T/\sqrt6$ have Laplacian
eigenvalues $1$ and $\kappa$, respectively. For $u=\sqrt d\,b_{12}$,
$t=1$, the exact solution is $R_1u=\sqrt{d/2}\,p$.
An allowed relative shifted-energy error is
\[
 e=\eta\sqrt{d/(1+\kappa)}\,r,
 \qquad \norm e_{L+I}=\eta\norm{R_1u}_{L+I}.
\]
Nevertheless
$\norm{L^{-\alpha/2}e}_2
=\eta\sqrt{d/(1+\kappa)}\,\kappa^{-\alpha/2}$ diverges as
$\kappa\downarrow0$. Both vectors are already mean zero.
Thus the standard shifted energy guarantee cannot simply be inserted
into a proof of relative factor accuracy. The feasible-flow construction
above supplies the missing control.

\begin{thebibliography}{9}
\bibitem{KOSZ}
J. A. Kelner, L. Orecchia, A. Sidford and Z. Allen-Zhu,
\emph{A Simple, Combinatorial Algorithm for Solving SDD Systems in
Nearly-Linear Time}, STOC 2013, pp. 911--920.
\url{https://arxiv.org/abs/1301.6628}.
Primary manuscript consulted:
\url{https://people.eecs.berkeley.edu/~brecht/cs294docs/week5/13.kelner.sdd.pdf}.
\end{thebibliography}
\end{document}
